% =================================================
% Spanish abstract
% =================================================
\cleardoublepage
\thispagestyle{empty}

\begin{otherlanguage}{spanish}

\begin{center}
    {\large\bfseries
    Compensación de la heterogeneidad de los rangos glucémicos en modelos
    de predicción de glucosa mediante balanceo de datos
    \par}

    \vspace{0.5cm}

    Diego Santiago Ortiz
\end{center}

\vspace{0.7cm}

\noindent
\textbf{Palabras clave}: diabetes tipo 1, T1D, monitorización continua de glucosa,
predicción de glucosa, desbalanceo de datos, balanceo de datos, SMOTER, SMOGN.

\vspace{0.7cm}

\noindent
\textbf{Resumen}

\vspace{0.3cm}

La predicción de la glucosa mediante el uso de modelos de Inteligencia Artificial puede facilitar una intervención temprana y segura ante episodios críticos de pacientes con Diabetes Tipo 1. Sin embargo, los modelos de regresión entrenados a partir de datos reales de Monitorización Continua de Glucosa tienden a favorecer los estados glucémicos más frecuentes, mientras muestran un rendimiento deficiente en regiones poco frecuentes pero clínicamente críticas, especialmente en hipoglucemia.

Este Trabajo de Fin de Grado estudia el efecto de balancear los diferentes rangos glucémicos de un dataset, sobre el rendimiento de modelos de predicción de glucosa para personas con Diabetes Tipo 1. Se utilizan 3 datasets de MCG: DiaTrend, REPLACE-BG, y T1DiabetesGranada. Se aplica balanceo de manera híbrida: submuestreo de los rangos sobrerrepresentados y sobremuestreo de los infrarrepresentados. Se exploran tres distribuciones: balanceo uniforme, una referencia normalizada y semi-balanceo. Cada una de estas distribuciones objetivo se combina con tres métodos de sobremuestreo: aleatorio, SMOTER y SMOGN. La arquitectura, el entrenamiento y los conjuntos de validación y test no se modifican. El rendimiento conseguido se evalúa estadísticamente usando tanto métricas de regresión como medidas clínicas derivadas de Clarke Error Grid.

Aplicando balanceo uniforme a todos los rangos se consiguen mejoras significativas en la proporción de predicciones que son clínicamente aceptables para la hipoglicemia severa, con aumentos de 63.75, 76.81, y 75.93 puntos porcentuales respectivamente para los 3 datasets. Aunque los errores de regresión aumentaron moderadamente a nivel global, las predicciones clínicamente aceptables permanecieron por encima del 99\%, sin que ninguna predicción adicional entrara en zonas peligrosas de la cuadrícula de Clarke. En resumen, este trabajo demuestra que, equilibrando la distribución de los rangos glucémicos, se puede mejorar sustancialmente la predicción en rangos minoritarios clínicamente críticos, manteniendo al mismo tiempo un elevado rendimiento clínico a nivel global.

\end{otherlanguage}

% =================================================
% English abstract
% =================================================
\cleardoublepage
\thispagestyle{empty}

\begin{center}
    {\large\bfseries
    Compensating for Glycemic Range Heterogeneity in Glucose Prediction
    Models via Data Balancing
    \par}

    \vspace{0.5cm}

    Diego Santiago Ortiz
\end{center}

\vspace{0.7cm}

\noindent
\textbf{Keywords}: Type 1 Diabetes, T1D, continuous glucose monitoring,
glucose prediction, data imbalance, data balancing, SMOTER, SMOGN.

\vspace{0.7cm}

\noindent
\textbf{Abstract}

\vspace{0.3cm}

Glucose prediction using AI models can support early and safe intervention of T1D critical episodes. However, regression models trained on real-world Continuous Glucose Monitoring data tend to favour the most frequent glycemic states while performing poorly in rare but clinically critical regions -- particularly hypoglycemia. This Bachelor’s Thesis investigates whether balancing the glycemic ranges proportions in training data can reduce this bias and improve the prediction of these critical events.

This Bachelor’s Thesis studies how balancing the different glycemic ranges affects glucose prediction models for T1D. It uses three CGM datasets: DiaTrend, REPLACE-BG, and T1DiabetesGranada. Balancing is approached hybridly: undersampling over-represented ranges and oversampling under-represented ones. Three distributions are explored: fully-balanced, normalized baseline, and semi-balancing. Each target distribution is combined with three oversampling methods: random resampling, SMOTER, and SMOGN. Predictive architecture, training pipeline, and validation and test sets remain unchanged. Performance is statistically assessed using both regression metrics and clinically meaningful measures derived from the Clarke Error Grid.

Full balancing all the glycemic ranges significantly increased the proportion of clinically acceptable severe-hypoglycemia predictions, with mean improvements of 63.75, 76.81, and 75.93 percentage points across the three datasets. Although convential regression metrics moderately increased globally, clinical acceptability remained above 99\%, with no predictions entering dangerous Clarke zones. Overall, the study demonstrates that modifying the glycemic distribution of the training data can substantially improve prediction in clinically critical minority ranges while preserving strong global clinical performance.

% =================================================
% Repository authorization
% =================================================
\cleardoublepage
\thispagestyle{empty}

\begin{otherlanguage}{spanish}

\noindent
\rule{\textwidth}{2pt}

\vspace{1cm}

Yo, \textbf{Diego Santiago Ortiz}, alumno del Grado en Ingeniería Informática
de la \textbf{Escuela Técnica Superior de Ingenierías Informática y de
Telecomunicación de la Universidad de Granada}, con DNI
\textbf{20889444Q}, autorizo la inclusión de una copia de este Trabajo Fin de
Grado en la biblioteca del centro para que pueda ser consultada por las
personas interesadas.

\vspace{6cm}

\noindent
Fdo.: Diego Santiago Ortiz

\vspace{2cm}

\begin{flushright}
Granada, a \textbf{[día]} de \textbf{[mes]} de 2026.
\end{flushright}

\end{otherlanguage}


% =================================================
% Supervisors' authorization
% =================================================
\cleardoublepage
\thispagestyle{empty}

\begin{otherlanguage}{spanish}

\noindent
\rule{\textwidth}{2pt}

\vspace{1cm}

D. \textbf{Oresti Baños Legrán}, profesor del Área de
\textbf{Ingeniería de Sistemas y Automática} del Departamento de
\textbf{Ingeniería de Sistemas y Automática} de la Universidad de Granada.

\vspace{0.5cm}

Dña. \textbf{Claudia Villalonga Palliser}, profesora del Área de
\textbf{Arquitectura y Tecnología de Computadores} del Departamento de
\textbf{Arquitectura y Tecnología de Computadores} de la Universidad de Granada.

\vspace{1cm}

\textbf{Informan:}

\vspace{0.5cm}

Que el presente trabajo, titulado
\textit{\textbf{Compensating for Glycemic Range Heterogeneity in Glucose
Prediction Models via Data Balancing}}, ha sido realizado bajo su supervisión
por \textbf{Diego Santiago Ortiz}, y autorizan su defensa ante el tribunal
correspondiente.

\vspace{0.5cm}

Y para que conste, expiden y firman el presente informe en Granada, a
\textbf{[día]} de \textbf{[mes]} de 2026.

\vspace{1cm}

\textbf{Los supervisores:}

\vspace{5cm}

\noindent
\begin{minipage}[t]{0.48\textwidth}
    \centering
    \textbf{Oresti Baños Legrán}
\end{minipage}
\hfill
\begin{minipage}[t]{0.48\textwidth}
    \centering
    \textbf{Claudia Villalonga Palliser}
\end{minipage}

\end{otherlanguage}


% =================================================
% Acknowledgements
% =================================================
\cleardoublepage

\chapter*{Agradecimientos}
\addcontentsline{toc}{chapter}{Acknowledgements}
\thispagestyle{empty}

Agradecimientos en español